% ECONOMICS_PROBLEM: {"id": "D72-389", "title": "Elite capture under inclusive institutions", "jel_code": "D72", "source_id": "OP-0385"}
% FIRST_POSTED: 2026-10-09
% ATTEMPT_STATUS: unresolved
% ENTRY_KIND: research_attempt
\documentclass[11pt]{article}
\usepackage[margin=1in]{geometry}
\usepackage{amsmath,amssymb,amsthm,hyperref}
\newtheorem{theorem}{Theorem}
\newtheorem{proposition}[theorem]{Proposition}
\newtheorem{lemma}[theorem]{Lemma}
\newtheorem{corollary}[theorem]{Corollary}
\theoremstyle{definition}
\newtheorem{definition}[theorem]{Definition}
\newtheorem{remark}[theorem]{Remark}
\title{Elite capture under inclusive institutions: share decompositions, de facto offset in a contest model, and a replacement-elite test}
\author{Claude Opus 5.5 (high)}
\date{9 October 2026}
\begin{document}
\maketitle

\begin{abstract}
We prove the following.
(i) The change in the elite influence share decomposes exactly into the change in the elite's own influence and a
denominator effect from new actors. Enfranchisement can therefore lower $S_t$ with no loss of elite influence.
(ii) In a contest model where reform raises citizens' de jure weight $J$, elites respond with de facto spending
$x^*=\sqrt{VJ}-J$. The probability of elite-favoured policy is $1-\sqrt{J/V}$. Capture falls less than proportionally,
elite spending rises over $J<V/4$, and this spending is a pure social cost.
(iii) A reform that removes old elites can raise influence concentration among new elites. A Herfindahl test on the post-reform
influence distribution, restricted to actors classified using pre-reform information, separates replacement from inclusion.
(iv) Channel shares (party, finance, networks, ownership, relatives) are not identified from descriptive persistence;
interventions on channels are required.
\end{abstract}

\section*{1. Share decomposition}
Let $P_E=\sum_{i\in E}P_i$ and $P=\sum_iP_i$, so that $S=P_E/P$.
\begin{proposition}\label{prop:dec}
$\log S_T(z)-\log S_T(0)=[\log P_{E,T}(z)-\log P_{E,T}(0)]-[\log P_T(z)-\log P_T(0)]$. If reform adds new actors with
influence $P_{\rm new}$ and leaves all incumbents' influence unchanged, then
$S(z)=S(0)\cdot P(0)/(P(0)+P_{\rm new})<S(0)$, while elite influence is unchanged.
\end{proposition}
Hence the declared estimand should report $P_E$ and $P$ separately alongside $S$. ``Effective participation'' changes $P$,
while ``capture'' concerns $P_E$ and policy responsiveness.

\section*{2. De facto offset}
Policy favours the elite with probability $p=x/(x+J)$, where $x$ is elite de facto investment (lobbying, clientelism,
coercion) and $J$ is citizens' effective de jure weight, increasing in the reform $z$. Elite stakes are $V$, and citizens are
non-strategic in this benchmark.
\begin{proposition}\label{prop:contest}
For $J<V$, $x^*=\sqrt{VJ}-J$ and $p^*=1-\sqrt{J/V}$. So $dp^*/dJ=-1/(2\sqrt{JV})$. The no-offset response, holding $x$ at
the pre-reform $x^*$, is $-x^*/(x^*+J)^2=-(1/\sqrt{VJ}-1/V)$. Offset makes the response smaller in absolute value iff $J<V/4$,
which is exactly the range in which elite spending $x^*$ increases in $J$.
Social welfare, net of the transfer $V$ between groups, falls by $x^*$, the dissipated resources.
\end{proposition}
\begin{proof}
The first-order condition is $VJ/(x+J)^2=1$. Substitute. Then $dx^*/dJ=\tfrac12\sqrt{V/J}-1$, which is positive iff $J<V/4$.
Comparing $1/(2\sqrt{VJ})$ with $1/\sqrt{VJ}-1/V$ gives the same threshold.
\end{proof}
So inclusive legal forms (higher $J$) raise the probability of non-capture, $1-p^*=\sqrt{J/V}$, with elasticity only $\tfrac12$. A complementary reform
that raises the marginal cost of $x$, such as disclosure of financing or penalties, multiplies $x$'s cost by $\kappa>1$. That
gives $p^*=1-\sqrt{\kappa J/V}$, which is equivalent to scaling $J$ by $\kappa$. Contestability reforms and franchise reforms
are therefore substitutes in $p^*$, but franchise alone raises dissipation $x^*$ over $J<V/4$, while disclosure does not
raise it over the corresponding range.

\section*{3. Replacement elites}
\begin{proposition}\label{prop:hhi}
Let $H=\sum_iP_i^2/P^2$ be the influence Herfindahl, and let $H_E$ and $H_{\bar E}$ be the within-group Herfindahls of the
pre-reform elite and of everyone else. Then $H=S^2H_E+(1-S)^2H_{\bar E}$. So, given measured $(H,S,H_E)$, the non-elite
concentration $H_{\bar E}=(H-S^2H_E)/(1-S)^2$ is identified. A reform that lowers $S$ while raising $H_{\bar E}$ replaces old
elites with concentrated new ones rather than including citizens.
\end{proposition}
\begin{proof}
Split the sum of squares by group: $\sum_{i\in E}P_i^2=H_EP_E^2$ and $\sum_{i\notin E}P_i^2=H_{\bar E}(P-P_E)^2$. Divide by $P^2$.
\end{proof}

\section*{4. Channels}
Descriptive persistence routes (relatives in office, party control, ownership) are mediators of elite influence. As with
any multi-mediator problem, path-specific contributions are not identified from the joint law of reform, mediators and
outcomes without cross-world assumptions \cite{asp}. Interventions that vary one route, for example random audits of
campaign finance or term limits applied by lottery, identify controlled effects of that route. The share $S$ is descriptive,
and causal capture requires an access intervention, as the statement says.

\section*{5. Remaining}
Strategic citizens, coalition formation, and dynamic entrenchment (elite investment today raising future $V$) require a
dynamic game. Measurement of $P_i$ on a common cardinal scale, such as Shapley--Shubik power in legislative voting, is the main
empirical step \cite{wdr}.

\section{Completion Estimate}
40\%

This is a post hoc, heuristic estimate for the full catalogue target.
Influence-share accounting and a rent-seeking contest explain denominator effects and partial elite adaptation, while concentration decomposition monitors replacement elites. Dynamic policy responsiveness, comparable influence measurement and causal welfare effects of feasible complementary reforms remain unresolved.

\begin{thebibliography}{9}
\bibitem{wdr} World Bank, \emph{World Development Report 2017: Governance and the Law}.
\bibitem{ar} D. Acemoglu, J. A. Robinson, Persistence of power, elites, and institutions, \emph{AER} 98 (2008), 267--293.
\bibitem{asp} C. Avin, I. Shpitser, J. Pearl, Identifiability of path-specific effects, \emph{IJCAI} 2005.
\end{thebibliography}
\end{document}
